\section{Limits of shrinking smooth potentials}\label{sec:hopf-defect-classification}

When the bounded mass property fails, \cref{cor:smooth-capacity-mass-defect} produces smooth potentials converging to zero whose Monge--Amp\`ere measures keep an excess mass in the limit. On the Hopf threefold we can determine exactly which excess measures occur: they are the lifts of arbitrary base measures by the fiber probability measures. The proof combines the intermediate currents of \cref{prop:shrinking-intermediate-currents} with averaging over the automorphism group.

On a general Hermitian manifold, the restrictions on the intermediate currents in \cref{prop:shrinking-intermediate-currents} are necessary; we do not know whether they are sufficient, or what the support and pluripolar behavior of such excess measures can be in general.

We keep the Hopf threefold $X$, its projection $\pi\colon X\to\mathbb P^2$, and the forms $\omega=\alpha+\beta$ of the previous sections. For $p\in\mathbb P^2$, let $h_p$ be the Haar probability measure on the elliptic fiber $\pi^{-1}(p)$, that is, the measure $\mathrm{d}A_E$ transported by the action of $E$. For a finite positive measure $\mu$ on the base, define its \emph{fiberwise lift} $\mathcal L\mu$ by
\[
 \int_X\varphi\,\mathrm{d}\mathcal L\mu
 =\int_{\mathbb P^2}\left(\int_{\pi^{-1}(p)}\varphi\,\mathrm{d}h_p\right)\mathrm{d}\mu(p),
 \qquad\varphi\in C(X).
\]
Since fiber integrals of continuous functions are continuous in $p$, the map $\mathcal L$ is weakly continuous. The normalization $\int_{\pi^{-1}(p)}\beta=2$ gives
\begin{equation}\label{eq:hopf-fiber-lift-normalization}
 \mathcal L\mu=\tfrac12\beta\wedge\pi^*\mu,
 \qquad \mathcal L(\omega_{\mathrm{FS}}^2)=\tfrac16\omega^3.
\end{equation}
Here $\pi^*\mu$ denotes the pullback of a top-degree current by the submersion $\pi$; locally it is the base measure with constant coefficients in the fiber directions.

For a Hermitian form $\Omega$ on $X$, put
\[
 M_\Omega(f)=\int_X(\Omega+\ddc f)^3,\qquad V_\Omega=\int_X\Omega^3
\]
for smooth real $f$, and let $b$ be a real number. A positive measure $\sigma$ on $X$ is called a \emph{shrinking defect} for $(\Omega,b)$\index{Mass defect!shrinking defect} if, for every sequence $0<\varepsilon_j<1$ tending to zero, there are smooth $U_j\in\PSH(X,\Omega)$ and numbers $0<r_j<\varepsilon_j$ such that
\begin{equation}\label{eq:hopf-shrinking-realization}
 \sup_XU_j=0,\qquad w_j=r_jU_j,\qquad
 M_\Omega(w_j)=b,\qquad
 (\Omega+\ddc w_j)^3\rightharpoonup\Omega^3+\sigma.
\end{equation}
Here weak convergence means convergence against every continuous function on $X$. We first classify the shrinking defects for the metric $\omega$, for which $V_\omega=6$; we then write $m=b-6>0$. Afterwards we pass to an arbitrary Hermitian form.

\subsection{Positive quadratic measures}

The mass formula of \cref{prop:hopf-mass} gives
\[
 M_\omega(\lambda f)-6=\lambda^2\bigl(M_\omega(f)-6\bigr)
\]
for every smooth real $f$ and real $\lambda$. To use this scalar identity in a weak limit, we also need a positive measure that represents the quadratic term.

\begin{lemma}\label{lem:hopf-positive-quadratic-measure}
Let $f$ be smooth, $0<t<1$, and $\ddc f\geq-t\omega$. Set $H=\ddc f+t\omega$, $d=M_\omega(f)-6$, and $\nu=3\omega\wedge H^2$. Then $d\geq0$ and
\[
 \begin{gathered}
 \int_X\omega^2\wedge H=6t,\qquad
 \int_X\omega\wedge H^2=d/3+6t^2,\qquad
 \int_XH^3=td+6t^3,\\
 \bigl\|(\omega+\ddc f)^3-\omega^3-\nu\bigr\|_{\mathrm{TV}}
 \leq36t+2td+24t^3.
 \end{gathered}
\]
\end{lemma}
\begin{proof}
The Gauduchon identity $\ddc\omega^2=0$ gives $\int_X\omega^2\wedge\ddc f=0$, and Stokes' formula gives $\int_X(\ddc f)^3=0$. Expanding the three integrals with these identities proves the asserted values, and \cref{prop:hopf-mass} gives $d\geq0$. The identity of measures
\[
 (\omega+\ddc f)^3
 =(1-t)^3\omega^3+3(1-t)^2\omega^2\wedge H
   +(1-t)\nu+H^3
\]
has four positive terms. After subtracting $\omega^3+\nu$, the total masses of the remaining terms bound the total variation by $18t+18t+t(d+18t^2)+td+6t^3$, as asserted.
\end{proof}

\subsection{Orbital mixtures}

Averaging a shrinking defect over automorphisms of $X$ produces new shrinking defects. The compact group $G=U(3)\times(\mathbb R/\mathbb Z)$ acts transitively on $X$ by the holomorphic maps $g[z]=[\mathrm{e}^tAz]$, $g=(A,t)$, and preserves $\omega$. Let $h_G$ be its Haar probability measure. For a probability measure $\rho$ on $G$ and a finite measure $\sigma$ on $X$, write
\[
 \mathcal A_\rho\sigma=\int_Gg^*\sigma\,\mathrm{d}\rho(g).
\]
Here the pullback of a measure is defined by $\int_X\varphi\,\mathrm{d}(g^*\sigma)=\int_X\varphi\circ g^{-1}\,\mathrm{d}\sigma$. In particular, transitivity gives
\[
 \int_Gg^*\zeta\,\mathrm{d}h_G(g)
 =\frac{\zeta(X)}6\omega^3
\]
for every finite measure $\zeta$ on $X$.

\begin{proposition}\label{prop:hopf-orbital-defects}
\index{Mass defect!orbital mixtures}
If $\sigma$ is a shrinking defect for $(\omega,6+m)$, then so is $\mathcal A_\rho\sigma$ for every probability measure $\rho$ on $G$.
\end{proposition}
\begin{proof}
Fix a sequence $\varepsilon_j$ as in the definition. Take a sequence $f_k=t_ku_k$ realizing $\sigma$ as in \eqref{eq:hopf-shrinking-realization}, so that $f_k$ has mass $6+m$, and put $H_k=\ddc f_k+t_k\omega$. By \cref{lem:hopf-positive-quadratic-measure}, the positive measures $\nu_k=3\omega\wedge H_k^2$ have masses $m+18t_k^2$ and converge weakly to $\sigma$.

The idea is to replace $f_k$ by a normalized sum of $N$ random translates. First suppose that $\rho$ has density at most $L$ relative to $h_G$. Choose $N$ independent random group elements $g_1,\ldots,g_N$ with law $\rho$, and set
\[
 W=N^{-1/2}\sum_{a=1}^Nf_k\circ g_a,\qquad
 \tau=t_k\sqrt N,\qquad
 K=\ddc W+\tau\omega=N^{-1/2}\sum_{a=1}^Ng_a^*H_k\geq0.
\]
The quadratic measure of $W$ is
\[
 3\omega\wedge K^2=D+C,\qquad
 D=N^{-1}\sum_{a=1}^Ng_a^*\nu_k,\qquad
 C=\frac6N\sum_{a<c}\omega\wedge g_a^*H_k\wedge g_c^*H_k\geq0.
\]
Averaging $u_k\circ g$ against $h_G$ gives a constant, hence $\int_Gg^*H_k\,\mathrm{d}h_G=t_k\omega$. Consequently $B=\int_Gg^*H_k\,\mathrm{d}\rho\leq Lt_k\omega$, and independence gives
\[
 \mathbb E[C(X)]=3(N-1)\int_X\omega\wedge B^2
 \leq18(N-1)L^2t_k^2.
\]
For any real $\varphi\in C(X)$ with $\|\varphi\|_\infty\leq1$, the variance of $D(\varphi)$ is at most $(m+18t_k^2)^2/N$.

For general $\rho$, convolve it with the normalized restrictions of $h_G$ to shrinking neighborhoods of the identity, which have positive Haar measure. This gives probability measures $\rho_j\rightharpoonup\rho$ with bounded densities $L_j$. Fix a countable dense set $(\varphi_l)$ in the real unit ball of $C(X)$ and numbers $0<\delta_j<\min(1,m/4)$ tending to zero. Choose $N_j$ large enough that
\[
 \frac{j(m+18)^2}{N_j\delta_j^2}<\frac14.
\]
Next choose $k_j$ large enough that $t_{k_j}<1$, $\tau_j=t_{k_j}\sqrt{N_j}<\min(\varepsilon_j/4,1/j,1/4)$, $18N_jt_{k_j}^2<m/4$, and the expected cross-term mass is less than $\delta_j/4$. Chebyshev's inequality for the first $j$ test functions and Markov's inequality for $C(X)$ show that some choice of $N_j$ group elements satisfies
\[
 |D_j(\varphi_l)-(\mathcal A_{\rho_j}\nu_{k_j})(\varphi_l)|<\delta_j
 \quad(l\leq j),\qquad C_j(X)<\delta_j.
\]
We may also arrange $k_j\to\infty$. For each fixed test function $\varphi$, the family $\{\varphi\circ g^{-1}:g\in G\}$ is compact in $C(X)$, so a finite uniform net and the bounded masses of $\nu_k$ show that $\int\varphi\circ g^{-1}\,\mathrm{d}\nu_k\to\int\varphi\circ g^{-1}\,\mathrm{d}\sigma$ uniformly in $g$. It follows that $\mathcal A_{\rho_j}\nu_{k_j}\rightharpoonup\mathcal A_\rho\sigma$. The dense test functions and the bounded masses now imply $D_j\rightharpoonup\mathcal A_\rho\sigma$, while $C_j\to0$ in total variation.

By the integral identity of \cref{lem:hopf-positive-quadratic-measure},
\[
 d_j=M_\omega(W_j)-6
 =m+C_j(X)-18(N_j-1)t_{k_j}^2\longrightarrow m,
 \qquad d_j>3m/4.
\]
Put $\lambda_j=\sqrt{m/d_j}$, so $\lambda_j\to1$ and $\lambda_j<2$. The average $a_j=N_j^{-1}\sum_a u_{k_j}\circ g_a$ is smooth and $\omega$-plurisubharmonic. With
\[
 U_j=a_j-\sup_Xa_j,\qquad r_j=\lambda_j\tau_j<\varepsilon_j,
\]
the potential $r_jU_j$ differs from $\lambda_jW_j$ by a constant and has exact mass $6+m$. Its positive quadratic measure is $\lambda_j^2(D_j+C_j)$. The total-variation estimate of \cref{lem:hopf-positive-quadratic-measure}, with parameter $r_j$, proves the required weak limit.
\end{proof}

Since $v_+(\omega)=+\infty$, \cref{cor:smooth-capacity-mass-defect} gives a sequence with some defect of mass $m$. Passing to a further subsequence meets any prescribed bounds $\varepsilon_j$ without changing the defect, so a shrinking defect of mass $m$ exists. Applying \cref{prop:hopf-orbital-defects} with $\rho=h_G$ shows that $(m/6)\omega^3$ is a shrinking defect. In particular, the excess mass can persist as a smooth limiting density.

\subsection{Fiber measures}

The intermediate currents force every shrinking defect to be a lift of a base measure, and the automorphisms of $X$ realize all such lifts.

\begin{lemma}\label{lem:hopf-defect-fiber-necessity}
Let $f_j=t_ju_j$, where $u_j\in\PSH(X,\omega)$ are smooth, $\sup_Xu_j=0$ and $0<t_j\to0$. Suppose that $M_\omega(f_j)=6+m$ and $(\omega+\ddc f_j)^3\rightharpoonup\omega^3+\sigma$. Then there is a unique positive measure $\mu$ on $\mathbb P^2$, of mass $m$, such that
\[
 \sigma=\mathcal L\mu,\qquad
 (\ddc f_j+t_j\omega)^2\rightharpoonup\tfrac16\pi^*\mu.
\]
\end{lemma}
\begin{proof}
Put $H_j=\ddc f_j+t_j\omega$. By \cref{prop:shrinking-intermediate-currents}, applied with $h=\omega$ and $n=3$, every convergent subsequence of $H_j^2$ has a positive closed limit $K$ with $\sigma=3\omega\wedge K$; moreover $\alpha\wedge K=0$, because $\alpha$ is closed and semipositive, so that $\alpha\wedge K$ is a positive measure with $\int_X\alpha\wedge K=0$. In holomorphic coordinates $(\xi_1,\xi_2,w)$ on a local trivialization of the elliptic bundle, the coefficient matrix of $K$ on the three exterior pairs is positive semidefinite as a matrix of measures. To see this directly, note that positivity of the matrices of $H_j^2$ passes to the limit when they are tested against constant vectors, and take Radon--Nikodym derivatives with respect to the diagonal trace. Since $\alpha$ is strictly positive in the horizontal directions, the identity $\alpha\wedge K=0$ makes both diagonal entries containing $\mathrm{d}w$ vanish. Positivity then makes their rows and columns vanish, so only the horizontal top-degree coefficient remains.

Since $K$ is closed, both real fiber derivatives of that coefficient vanish in the sense of distributions. Testing against a compactly supported base function times a nonconstant Fourier mode on $E=\mathbb C/(\mathbb Z+2\pi\mathrm{i}\mathbb Z)$ shows that every such mode vanishes. By density of trigonometric polynomials in $C(E)$, the coefficient is a base measure times the constant fiber area. Changes of bundle coordinates translate $w$ by a holomorphic function of the base and preserve that area. Thus $K=\pi^*\tau$ for a positive measure $\tau$ on $\mathbb P^2$.

Since $\alpha\wedge K=0$, \eqref{eq:hopf-fiber-lift-normalization} gives $\sigma=3\beta\wedge\pi^*\tau=\mathcal L(6\tau)$. Therefore $\mu=6\tau=\pi_*\sigma$ is uniquely determined by $\sigma$, and it has mass $m$. In particular, every subsequential limit of $H_j^2$ is the same, which proves convergence of the entire sequence.
\end{proof}

\begin{lemma}\label{lem:hopf-fiber-defect-realization}
For every $p\in\mathbb P^2$, the measure $mh_p$ is a shrinking defect for $(\omega,6+m)$.
\end{lemma}
\begin{proof}
Start with smooth $f_j=t_ju_j$ realizing the shrinking defect $(m/6)\omega^3=\mathcal L(m\,\omega_{\FS}^2)$, and put $H_j=\ddc f_j+t_j\omega$. By \cref{lem:hopf-defect-fiber-necessity}, $H_j^2\rightharpoonup(m/6)\alpha^2$, while $H_j$ has trace mass $\int_X\omega^2\wedge H_j=6t_j$ by \cref{lem:hopf-positive-quadratic-measure}. Fix $A\in GL(3,\mathbb C)$, regarded as an automorphism of $X$, and let $a$ be the induced projective transformation of $\mathbb P^2$. Choose $C_A\geq1$ with $A^*\omega\leq C_A\omega$, and put
\[
 R_A=C_A\omega-A^*\omega\geq0,\qquad
 H_j^A=\ddc(f_j\circ A)+C_At_j\omega=A^*H_j+t_jR_A.
\]
For this fixed $A$, the trace mass of $A^*H_j$ is $O_A(t_j)$, and the positive difference $(H_j^A)^2-A^*(H_j^2)$ has trace mass $O_A(t_j^2)$. Since $\alpha=\pi^*\omega_{\FS}$, it follows that
\[
 (H_j^A)^2\rightharpoonup\frac m6\pi^*(a^*\omega_{\mathrm{FS}}^2),
 \qquad 3\omega\wedge(H_j^A)^2\rightharpoonup\mathcal L(m a^*\omega_{\mathrm{FS}}^2).
\]
The limiting measure has mass $m$, so the integral identity of \cref{lem:hopf-positive-quadratic-measure}, applied with the parameter $C_At_j$, gives $d_j^A=M_\omega(f_j\circ A)-6\to m$. For large $j$, multiply $f_j\circ A$ by $\sqrt{m/d_j^A}\to1$. The mass is then exactly $6+m$, and the weak limit is unchanged. The underlying normalized potential is $u_j\circ A/C_A\in\PSH(X,\omega)$, with parameter $C_At_j\sqrt{m/d_j^A}\to0$. The total-variation estimate of \cref{lem:hopf-positive-quadratic-measure} therefore shows that $\mathcal L(m\,a^*\omega_{\FS}^2)$ is a shrinking defect; any prescribed bounds $\varepsilon_j$ are met by passing to a sufficiently late subsequence.

Now take $A_l=\operatorname{diag}(R_l^{-1},1,1)$, $R_l\to+\infty$. Since
\[
 a_l^*\omega_{\mathrm{FS}}^2=(a_l^{-1})_*\omega_{\mathrm{FS}}^2,
\]
and $a_l^{-1}[q_0:q_1:q_2]=[R_lq_0:q_1:q_2]\to[1:0:0]$ outside the hyperplane $\{q_0=0\}$, which has zero volume, dominated convergence gives $a_l^*\omega_{\mathrm{FS}}^2\rightharpoonup\delta_{[1:0:0]}$. By weak continuity of $\mathcal L$, the shrinking defects $\mathcal L(m\,a_l^*\omega_{\FS}^2)$ converge to $mh_{[1:0:0]}$. For each $l$, choose one corrected potential for $A_l$, sufficiently late in its sequence that its parameter is less than $\varepsilon_l$ and its measure differs from the corresponding limit by less than $1/l$ on the first $l$ functions of a fixed dense sequence of test functions. This diagonal sequence has the exact mass $6+m$ and the limit $\omega^3+mh_{[1:0:0]}$. Unitary automorphisms move this fiber to any prescribed fiber.
\end{proof}

\Needspace{9\baselineskip}
\begin{theorem}\label{thm:hopf-standard-defect-classification}
\index{Hopf manifold!shrinking defects}\index{Mass defect!fiberwise classification}
For $b>6$, the shrinking defects for $(\omega,b)$ are exactly $\mathcal L\mu$, where $\mu$ ranges over all positive measures on $\mathbb P^2$ of mass $b-6$. Each such defect has support $\pi^{-1}(\operatorname{supp}\mu)$ and has no point atoms. It charges no pluripolar sets if and only if $\mu$ charges no pluripolar sets on $\mathbb P^2$.
\end{theorem}
\begin{proof}
Put $m=b-6$. Necessity follows from \cref{lem:hopf-defect-fiber-necessity}. Finite convex combinations of the measures $mh_p$ are shrinking defects by \cref{lem:hopf-fiber-defect-realization,prop:hopf-orbital-defects}, the latter applied to an atomic measure $\rho$ supported on unitary maps that move a fixed fiber to the desired fibers. Approximate an arbitrary base measure of mass $m$ by finite atomic measures of the same mass, using finite partitions of diameter tending to zero. Weak continuity of $\mathcal L$ and the diagonal choice on dense test functions used in the preceding proof give a realizing sequence for every prescribed sequence $\varepsilon_j$.

Every nonempty open subset of a fiber has positive Haar measure. Local trivializations therefore show that $\mathcal L\mu$ has support $\pi^{-1}(\operatorname{supp}\mu)$, and the fiber measures have no point atoms. If $B$ is a pluripolar subset of the base, local plurisubharmonic functions whose pole sets contain $B$ pull back by the submersion to such functions on $X$. Hence $\pi^{-1}(B)$ is pluripolar, and $\mathcal L\mu(\pi^{-1}(B))=\mu(B)$; thus $\mathcal L\mu$ charges a pluripolar set whenever $\mu$ does.

Conversely, suppose that $\mu$ charges no pluripolar subset of the base. Cover a pluripolar set in $X$ by countably many local product charts, each carrying a plurisubharmonic function $q(\xi,w)\not\equiv-\infty$ whose pole set contains the part of the set in that chart. For a relatively compact disk $D$ in the fiber direction, the function
\[
 F(\xi)=\frac1{|D|}\int_Dq(\xi,w)\,\mathrm{d}A(w)
\]
is plurisubharmonic and not identically $-\infty$ on a smaller base chart. Indeed, local integrability of $q$ and Fubini's theorem show that $F$ is locally integrable, and in particular finite almost everywhere; a common upper bound and the reverse Fatou lemma give upper semicontinuity; and integrating the submean inequalities in the base directions gives the submean property. The restriction of $q$ to a fiber is either subharmonic with a pole set of zero area or identically $-\infty$. The fibers of the second kind lie over the pluripolar set $\{F=-\infty\}$, which has zero $\mu$-measure. Integrating the Haar measures of the fibers against $\mu$ therefore shows that each local pole set has zero $\mathcal L\mu$-measure, and the countable cover proves the assertion for this pluripolar set. For arbitrary pluripolar subsets, the same argument applies to Borel pole sets containing them, with outer measure. Similarly, for the first implication, a Borel $\mathcal L\mu$-null set containing $\pi^{-1}(B)$ has a measurable fiber-mass function equal to one on $B$, which gives the corresponding outer-measure statement.
\end{proof}

For every nonempty closed set $F\subseteq\mathbb P^2$, a measure $m\sum_{l\geq1}2^{-l}\delta_{p_l}$ carried by a countable dense subset $\{p_l\}$ of $F$ has support $F$. Thus every closed set of the form $\pi^{-1}(F)$ is the support of a shrinking defect. The fibers themselves are pluripolar in a global sense. For the fiber over $[1:0:0]$, the function
\[
 \psi=\log\frac{|z_1|^2+|z_2|^2}{|z|^2}
\]
satisfies $\ddc\psi\geq-\alpha$, hence $\omega+\ddc\psi\geq\beta\geq0$. Thus $\psi$ is $\omega$-plurisubharmonic, and its pole set is exactly that fiber; unitary maps give the other fibers. By \cref{thm:hopf-standard-defect-classification}, for $0\leq\lambda\leq1$ and $p\in\mathbb P^2$, the measure
\[
 m\lambda h_p+\frac{m(1-\lambda)}6\omega^3
\]
is a shrinking defect. It gives mass $m\lambda$ to the pluripolar fiber $\pi^{-1}(p)$, while every pluripolar set has mass at most $m\lambda$. For $0<\lambda<1$ its support is all of $X$, despite this positive mass on a pluripolar set.

\subsection{Changing the Hermitian metric}

For a general Hermitian form, the fiber Haar measures are replaced by the normalized area measures of the form. Let $\Omega$ be any Hermitian form on the Hopf threefold $X$. Its restriction to each elliptic fiber has positive area
\[
 A_\Omega(p)=\int_{\pi^{-1}(p)}\Omega.
\]
These areas depend smoothly on $p$ and are bounded above and below by positive constants. Define the probability measures $h_p^\Omega=\Omega|_{\pi^{-1}(p)}/A_\Omega(p)$ on the fibers, and define the lift $\mathcal L_\Omega\mu$ of a finite positive measure $\mu$ on $\mathbb P^2$ by the formula defining $\mathcal L\mu$, with $h_p^\Omega$ in place of $h_p$.

\begin{theorem}\label{thm:hopf-hermitian-defect-classification}
For every Hermitian form $\Omega$ on $X$ and every $b>V_\Omega$, the shrinking defects for $(\Omega,b)$ are exactly
\[
 \mathcal L_\Omega\mu,\qquad
 \mu\geq0,\qquad \mu(\mathbb P^2)=b-V_\Omega.
\]
Their supports are $\pi^{-1}(\operatorname{supp}\mu)$, they have no point atoms, and they charge no pluripolar sets exactly when $\mu$ charges no base pluripolar sets. Moreover, $v_+(\Omega)=+\infty$.
\end{theorem}
\begin{proof}
Put $m=b-V_\Omega>0$. For necessity, let $\sigma$ be a shrinking defect for $(\Omega,b)$. \Cref{prop:shrinking-intermediate-currents}, applied with $h=\Omega$, gives a positive closed limit $K$ of the squares of the intermediate forms with $\sigma=3\Omega\wedge K$, and $\alpha\wedge K=0$ as in the proof of \cref{lem:hopf-defect-fiber-necessity}. The local coefficient and Fourier argument of that proof gives $K=\pi^*\tau$ for a positive measure $\tau$ on $\mathbb P^2$. Consequently
\[
 \sigma=\mathcal L_\Omega\mu,\qquad
 \mu=3A_\Omega\tau=\pi_*\sigma,
\]
and $\mu$ has mass $m$.

Conversely, fix such a $\mu$ and set $\tau=\mu/(3A_\Omega)$. By \cref{thm:hopf-standard-defect-classification}, choose a sequence $f_j=t_ju_j$ realizing the shrinking defect $\mathcal L(6\tau)$ for $\omega$. The positive forms $H_j=\ddc f_j+t_j\omega$ have trace mass $O(t_j)$, and $H_j^2\rightharpoonup\pi^*\tau$ by \cref{lem:hopf-defect-fiber-necessity}. Choose $C\geq1$ with $\omega\leq C\Omega$, and put
\[
 R=C\Omega-\omega\geq0,\qquad
 G_j=\ddc f_j+Ct_j\Omega=H_j+t_jR\geq0.
\]
The trace mass of $G_j$ is $O(t_j)$, and $G_j^2-H_j^2$ has trace mass $O(t_j^2)$. Thus $G_j^2\rightharpoonup\pi^*\tau$ and $3\Omega\wedge G_j^2\rightharpoonup\mathcal L_\Omega\mu$.

Unlike $\omega$, the form $\Omega$ need not be Gauduchon, so the mass polynomial may have a linear term. Define
\[
 \ell_j=3\int_X\Omega^2\wedge\ddc f_j=O(t_j),\qquad
 q_j=3\int_X\Omega\wedge(\ddc f_j)^2\longrightarrow m.
\]
These estimates follow by expanding $\ddc f_j=G_j-Ct_j\Omega$ and using the preceding trace bounds. Stokes' formula removes the cubic integral, so
\[
 M_\Omega(\lambda f_j)=V_\Omega+\lambda\ell_j+\lambda^2q_j.
\]
For large $j$, $q_j>0$. The positive root
\begin{equation}\label{eq:hopf-hermitian-mass-correction}
 \lambda_j=\frac{-\ell_j+\sqrt{\ell_j^2+4mq_j}}{2q_j}
 \longrightarrow1
\end{equation}
gives $M_\Omega(\lambda_jf_j)=b$. The normalized function $U_j=u_j/C$ is $\Omega$-plurisubharmonic, and $w_j=\lambda_j f_j=r_jU_j$ with $r_j=\lambda_jCt_j\to0$.

To identify the full weak limit, first note that expanding $G_j=\ddc f_j+Ct_j\Omega$ gives
\[
 \int_XG_j^3=Ct_jq_j+(Ct_j)^2\ell_j+(Ct_j)^3V_\Omega=O(t_j).
\]
Hence $G_j^3\to0$ in total variation. The positive expansion
\[
 (\Omega+\ddc w_j)^3
 =(1-r_j)^3\Omega^3
  +3(1-r_j)^2\lambda_j\Omega^2\wedge G_j
  +3(1-r_j)\lambda_j^2\Omega\wedge G_j^2
  +\lambda_j^3G_j^3
\]
converges weakly to $\Omega^3+\mathcal L_\Omega\mu$, since the second and fourth terms tend to zero in total variation. Given prescribed bounds $\varepsilon_j$, pass to a sufficiently late subsequence with $r_j<\varepsilon_j$ at every index. This preserves the normalization, the exact mass and the weak limit.

The fiber probabilities $h_p^\Omega$ have smooth densities relative to $h_p$ that are bounded above and below by positive constants, uniformly in $p$. Thus $\mathcal L_\Omega\mu$ and $\mathcal L\mu$ are bounded by fixed positive multiples of each other, and the assertions about support, point atoms and pluripolar sets follow from \cref{thm:hopf-standard-defect-classification}. The global functions with fiber poles also remain available: if $\Omega\geq c\omega$, then $c\psi\in\PSH(X,\Omega)$.

Finally, for smooth $u\in\PSH(X,\omega)$ and $c>0$ with $\Omega\geq c\omega$,
\[
 \Omega+\ddc(cu)\geq c(\omega+\ddc u),\qquad
 M_\Omega(cu)\geq c^3M_\omega(u).
\]
Since $v_+(\omega)=+\infty$ by \cref{cor:hopf-lower-mass} and the supremum may be computed over smooth potentials (\cref{prop:bmp-metric}), it follows that $v_+(\Omega)=+\infty$.
\end{proof}
